\section{Primitive projections in the Bost--Connes system}
\label{sec:thermal}

\subsection{The equilibrium identity}

The Bost--Connes arithmetic system was introduced in \cite{BC}.
We use its multiplicative isometries $v_n$ and the range projections
$e_n=v_nv_n^*$, with
\begin{equation}
 v_n^*v_n=I,\qquad v_mv_n=v_{mn},\qquad
 e_me_n=e_{\operatorname{lcm}(m,n)},\qquad
 \sigma_t(v_n)=n^{\ii t}v_n.
 \label{eq:bc_relations}
\end{equation}
In the standard integer representation $v_n|k\rangle=|nk\rangle$,
the projection $e_n$ detects divisibility by $n$, and the Hamiltonian
is $\mathcal H|k\rangle=(\log k)|k\rangle$.
This representation has a normalized Gibbs trace only for $\beta>1$.
The full system nevertheless has positive KMS states for
$0<\beta\le1$; their existence and local measure structure are
established inputs, for which we use \cite[Introduction and Section~2]{Neshveyev}.

For analytic elements the KMS condition reads
$\varphi_\beta(ab)=\varphi_\beta(b\sigma_{\ii\beta}(a))$.
Apply it to $a=v_n$ and $b=v_n^*$:
\begin{equation}
 \varphi_\beta(e_n)
 =\varphi_\beta(v_n^*\sigma_{\ii\beta}(v_n))=n^{-\beta}.
 \label{eq:divisibility_probability}
\end{equation}
No analytic continuation of a divergent trace is used in this identity.

\begin{proposition}[Primitive thermal coefficient identity]
For $n\ge1$ define the temperature-independent projection
\begin{equation}
 P_n=\prod_{\ell\mid n}(I-e_\ell),\qquad P_1=I.
 \label{eq:primitive_projection}
\end{equation}
For a KMS state as above,
\begin{equation}
 \varphi_\beta(P_n)=\prod_{\ell\mid n}(1-\ell^{-\beta}).
 \label{eq:primitive_probability}
\end{equation}
Consequently the centered expectations at $\beta=2\omega$ are exactly
the coefficients in \eqref{eq:coefficients}.
\end{proposition}

\begin{proof}
The divisibility projections commute. Expanding the finite product in
\eqref{eq:primitive_projection} and using the lcm relation gives
\begin{equation}
 P_n=\sum_{d\mid\operatorname{rad}(n)}\mu(d)e_d,
\end{equation}
where $\operatorname{rad}(n)$ is the product of the distinct prime
divisors, $\mu$ is the M\"obius function and $e_1=I$.
Taking expectations and applying \eqref{eq:divisibility_probability}
gives the product in \eqref{eq:primitive_probability}. Multiplication
by $n^{(\beta-1)/2}$ gives \eqref{eq:coefficients}. Thus the result
does not require a separate probabilistic independence assumption.
\end{proof}

At a repeated prime, $P_{\ell^r}=P_\ell$: primitiveness is the same
event. Its contribution is
\begin{equation}
 c_{\ell^r}(\omega)
 =(1-\ell^{-2\omega})\ell^{r(\omega-1/2)}.
\end{equation}
The repetition dependence comes from centering; one has not introduced
a new primitive event for each $r$. For coprime integers the projections
combine by their distinct prime sets, recovering the required
multiplicativity. Hence the first-prime tangent
\eqref{eq:first_prime} and the full prime source
\eqref{eq:prime_generator} follow from the same coefficient law.

There is also a direct counting interpretation at integer exponent.
Inclusion-exclusion shows that
\begin{equation}
 J_d(n)=n^d\prod_{\ell\mid n}(1-\ell^{-d})
\end{equation}
counts tuples $(x_1,\ldots,x_d)$ modulo $n$ with
$\gcd(n,x_1,\ldots,x_d)=1$. For real $\beta>0$ the analogous product
defines $J_\beta(n)$, and $c_n(\omega)=J_{2\omega}(n)/n^{\omega+1/2}$.
The thermal construction supplies a real-parameter probability law;
it does not posit a noninteger number of discrete coordinates.

\subsection{An explicit positive model at finitely many places}

Let $S$ be a finite set of primes. On
$\bigotimes_{\ell\in S}\ell^2(\mathbb N_0)$, with occupation operators
$N_\ell$, take the self-adjoint Hamiltonian
\begin{equation}
 \mathcal H_S=\sum_{\ell\in S}(\log\ell)N_\ell.
 \label{eq:finite_hamiltonian}
\end{equation}
For every $\beta>0$ its Gibbs state is trace class, with
\begin{equation}
 \mathcal Z_S(\beta)=\prod_{\ell\in S}(1-\ell^{-\beta})^{-1},
 \qquad
 \Pr_\beta(N_\ell=m)=(1-\ell^{-\beta})\ell^{-m\beta}.
 \label{eq:geometric_law}
\end{equation}
The primitive projection selects occupation zero at each prime dividing
$n$, provided those primes belong to $S$. It gives
\eqref{eq:primitive_probability} directly. This finite-place realization
of the relevant diagonal observables is explicit and positive throughout
the temperature range of interest. It is not asserted to encode every
observable of the full arithmetic algebra.

As a small illustration, the centered coefficients are
\begin{center}
\small
\begin{tabular}{@{}rrrr@{}}
\toprule
$n$ & $\omega=1/4$ & $\omega=0.4$ & $\omega=1/2$\\
\midrule
2 & 0.246292858 & 0.397146260 & $1/2$\\
3 & 0.321144348 & 0.523917402 & $2/3$\\
4 & 0.207106781 & 0.370550563 & $1/2$\\
6 & 0.079095559 & 0.208071837 & $1/3$\\
\bottomrule
\end{tabular}
\end{center}
The entries are samples of the exact identity, not a fit defining the
state. In the infinite system for $\beta\le1$, the product partition
function diverges. The positive KMS state must then be used in its
established algebraic sense, not replaced by
$e^{-\beta\mathcal H}/\zeta(\beta)$.

\subsection{What the centering and the time variable still require}

For $0<\beta\le1$, the centered observable
\begin{equation}
 B_{n,\beta}=n^{(\beta-1)/2}P_n
\end{equation}
is a positive contraction and
$\varphi_\beta(B_{n,\beta})=c_n(\beta/2)$. However, this observable
depends on temperature. The equality alone does not derive all
amplitudes from a temperature change with a fixed physical detector.

The dynamics gives a structured analytic-time expression:
\begin{equation}
 c_n(\beta/2)=n^{-1/2}\varphi_\beta
 \left(P_n v_n^*\sigma_{-\ii\beta/2}(v_n)\right),\qquad
 \sigma_{-\ii\beta/2}(v_n)=n^{\beta/2}v_n.
 \label{eq:thermal_halfstep}
\end{equation}
It identifies a thermal half-step and the fixed factor $n^{-1/2}$.
Imaginary-time continuation is not unitary real-time propagation.
The physical input/output normalization that would implement
\eqref{eq:thermal_halfstep} is still to be derived.

There is a second, independent distinction. Since
$\sigma_t(P_n)=P_n$, a native real-time correlation of two such
projections is constant:
\begin{equation}
 \varphi_\beta(P_m\sigma_t(P_n))=\varphi_\beta(P_mP_n).
\end{equation}
The number $\log n$ in $\sigma_t(v_n)=n^{\ii t}v_n$ is an energy
difference or dilation label. It has not thereby become the causal
delay $\log n$ in \eqref{eq:coefficients}. That conversion is a
requirement on a coupling to propagating signals, not a consequence
of equilibrium positivity.
